\chapter{Dépannage}
\label{ch:troubleshooting}

Quand quelque chose tourne mal, commencez par les logs et le tableau
de bord \texttt{/admin > Health}. Ce chapitre catalogue les
problèmes que les utilisateurs rencontrent vraiment, par ordre de
fréquence approximatif.

\section{Lire les logs}

\begin{lstlisting}
docker compose logs --since 10m mybibli
docker compose logs -f mybibli
\end{lstlisting}

Les logs sortent en JSON structuré. Les champs les plus utiles sont
\texttt{level}, \texttt{target} (quel module a logué cette ligne),
\texttt{message}, et toute paire clé=valeur contextuelle que la
ligne porte (\texttt{user\_id}, \texttt{volume\_id}, \dots).

Pour filtrer :

\begin{lstlisting}
docker compose logs --since 10m mybibli | grep -i error
\end{lstlisting}

\section{L'application ne démarre pas}

\subsection*{Symptôme : le conteneur s'arrête immédiatement}

Lancez \texttt{docker compose logs mybibli} et cherchez l'une des
lignes :

\begin{itemize}
  \item \texttt{missing required environment variable:
        DATABASE\_URL}\index{dépannage!DATABASE\_URL} — le fichier
        \texttt{.env} est absent ou la variable n'est pas
        positionnée. Voir section~\ref{sec:env-vars}.
  \item \texttt{Failed to create database pool} — l'hôte / port /
        identifiants dans \texttt{DATABASE\_URL} sont faux, ou la
        base n'est pas joignable. Depuis l'hôte, essayez
        \texttt{docker compose exec mybibli sh -c 'apk add curl;
        curl <db-host>:<db-port>'} pour confirmer la joignabilité.
  \item \texttt{Failed to run database migrations} — généralement
        un souci de permission. Confirmez que l'utilisateur DB a
        \texttt{ALL PRIVILEGES} sur la base (voir
        section~\ref{sec:install-external-db}).
  \item \texttt{Address already in use} — le port hôte publié est
        pris. Changez \texttt{HOST\_PORT} dans \texttt{.env}.
\end{itemize}

\subsection*{Symptôme : chaque page renvoie 404}

Le conteneur est levé mais les routes manquent. Cause la plus
fréquente : vous avez tiré un tag d'image qui n'existe pas (faute
de frappe) et Docker est retombé sur un vieux tag obsolète.
\texttt{docker compose pull} puis redémarrez.

\section{Impossible de se connecter}

\subsection*{Symptôme : « Invalid credentials » sur une installation 1.0.x neuve}

Les identifiants seedés sont \texttt{admin/admin}. S'ils ne marchent
pas non plus :

\begin{itemize}
  \item la migration de base qui les crée a échoué silencieusement
        — cherchez les erreurs de migration dans les logs ;
  \item l'admin précédent a déjà tourné le mot de passe (vérifiez
        avec qui a installé mybibli) ;
  \item vous regardez une installation 1.1.0+, où les identifiants
        seedés sont bien plantés par la migration mais immédiatement
        soft-deleted par la garde issue-\#173 au démarrage —
        utilisez plutôt l'assistant à \texttt{/setup}.
\end{itemize}

\section{Utilisateurs résiduels dans la table \texttt{users} après l'assistant}
\label{sec:troubleshoot-seeded-users}

\subsection*{Symptôme : quatre utilisateurs en base, un seul créé dans l'assistant}

Si vous inspectez la table \texttt{users} directement (via
phpMyAdmin ou le CLI MariaDB) juste après avoir terminé
\texttt{/setup}, vous verrez \emph{quatre} lignes alors que vous
n'avez saisi qu'un seul nom d'utilisateur et un seul mot de passe.
C'est intentionnel. Trois de ces lignes ne peuvent pas servir à se
connecter.

\begin{longtable}{@{} l l c l @{}}
\toprule
\textbf{username} & \textbf{role} & \textbf{active} & \textbf{deleted\_at} \\
\midrule
\endhead
\texttt{admin}     & \texttt{admin}     & 1 & rempli, $\approx$ heure du premier boot \\
\texttt{librarian} & \texttt{librarian} & 1 & rempli, $\approx$ heure du premier boot \\
\texttt{SYSTEM}    & \texttt{system}    & 0 & NULL \\
\texttt{<votre-username>} & \texttt{admin} & 1 & NULL \\
\bottomrule
\end{longtable}

\begin{itemize}
  \item \texttt{admin} et \texttt{librarian} sont plantés par les
        migrations de seed (qui tournent à chaque installation neuve
        pour le flux dev / E2E) puis immédiatement
        \emph{soft-deleted} par la garde issue-\#173 au démarrage.
        Leur \texttt{deleted\_at} est rempli, donc la requête de
        login (qui filtre \texttt{deleted\_at IS NULL}) les refuse.
  \item \texttt{SYSTEM}\index{utilisateur SYSTEM} est un compte
        utilitaire non connectable, utilisé pour attribuer les
        entrées du journal d'audit produites par la tâche de fond
        de purge automatique à 30 jours. Son rôle est \texttt{system}
        (hors de l'énum de login \texttt{admin/librarian}), son
        \texttt{password\_hash} est la chaîne littérale
        \texttt{NO\_LOGIN\_SYSTEM\_USER\_HASH\_NOT\_A\_VALID\_ARGON2\_STRING},
        et son flag \texttt{active} vaut \texttt{0}.
\end{itemize}

\begin{tip}
Pour vérifier que les lignes soft-deleted sont inertes, essayez de
vous connecter avec \texttt{admin/admin} : la tentative est
refusée. Le panneau \texttt{/admin > Utilisateurs} filtre la ligne
\texttt{SYSTEM} ainsi que toute ligne avec \texttt{deleted\_at}
rempli, donc l'interface admin n'affiche que le compte que vous
avez créé dans l'assistant.
\end{tip}

Les deux lignes soft-deleted sont physiquement supprimées de la
table par le planificateur de purge automatique à 30 jours —
attendez-vous à les voir disparaître environ un mois après la date
de boot inscrite dans \texttt{deleted\_at}.

\subsection*{Symptôme : la connexion redirige vers \texttt{/login}}

Le cookie de session est rejeté par le navigateur. Causes :

\begin{itemize}
  \item \texttt{MYBIBLI\_COOKIE\_SECURE=true} sur un déploiement
        plain-HTTP — le navigateur refuse d'accepter un cookie
        Secure en HTTP. Positionnez la variable à \texttt{false}
        ou frontez l'instance en HTTPS.
  \item Cookies bloqués par politique du navigateur (fenêtre
        privée avec blocage strict des cookies tiers, etc.).
        Testez dans un profil par défaut.
\end{itemize}

\subsection*{Symptôme : verrouillé dehors — dernier admin désactivé}

La garde du dernier admin actif devrait empêcher cela, mais si un
admin précédent a été désactivé puis supprimé définitivement, la
récupération demande un accès SQL direct :

\begin{lstlisting}
docker compose exec db mariadb -u root -p"$MYSQL_ROOT_PASSWORD" mybibli \
    -e "UPDATE users SET deleted_at=NULL, active=1 WHERE username='votre-admin';"
\end{lstlisting}

Ensuite connectez-vous normalement.

\section{Les métadonnées ne se résolvent jamais}

\subsection*{Symptôme : le titre squelette reste non résolu}

Ouvrez \texttt{/admin > Health}. Regardez l'état de chaque
fournisseur. Une croix rouge veut généralement dire :

\begin{itemize}
  \item le fournisseur est vraiment down (rare pour BnF / Open
        Library ; plus fréquent pour les services plus petits) ;
  \item la clé d'API est absente ou invalide — ouvrez
        \texttt{/admin > System > Metadata Providers} et re-saisissez
        la clé ;
  \item l'hôte ne peut pas joindre l'upstream. Confirmez avec
        \texttt{docker compose exec mybibli sh -c "wget -q
        https://www.googleapis.com/books/v1 -O -"} ou équivalent.
\end{itemize}

\subsection*{Symptôme : seuls les ISBN français se résolvent}

Vous avez désactivé (ou jamais positionné) les clés Google Books /
Open Library. BnF est excellent pour les titres édités en France
mais limité pour les titres en langues étrangères. Ajoutez au
moins Open Library (sans clé requise).

\section{La recherche renvoie trop peu de résultats}

\subsection*{Symptôme : des titres connus ne remontent pas}

\begin{itemize}
  \item Confirmez que le titre n'est pas soft-deleted (vérifiez
        \texttt{/admin > Trash}). La recherche cache les éléments
        soft-deleted.
  \item Confirmez que les contributeurs du titre sont orthographiés
        comme vous avez cherché. La recherche mybibli est floue
        mais pas magique ; « Camus » trouvera « Albert Camus »
        mais « Camu » peut ne rien donner.
\end{itemize}

\subsection*{Symptôme : la recherche est lente}

Pour les collections au-delà de 10\,000 titres, les requêtes qui
combinent contributeur et texte libre peuvent prendre une
seconde ou deux. C'est une limitation connue ; des index prévus
combleront l'écart.

\section{Images de couverture manquantes}

\subsection*{Symptôme : chaque titre affiche le placeholder}\label{sec:troubleshoot-couvertures-manquantes}

Le répertoire des couvertures est vide ou illisible. Confirmez :

\begin{itemize}
  \item que \texttt{docker-compose.yml} déclare le volume nommé
        \texttt{mybibli\_covers} (ou un bind mount sur
        \texttt{COVERS\_HOST\_PATH}) mappé sur \texttt{/app/covers}.
        Les installs antérieures à 1.1.4 ne déclaraient PAS ce
        volume — voir la sous-section suivante ;
  \item que la récupération de métadonnées a effectivement renvoyé
        une URL de couverture (regardez la page du titre — si l'URL
        manque, le fournisseur ne l'a pas fournie).
\end{itemize}

\subsection*{Symptôme : les couvertures ont disparu après mise à jour de l'image}\label{sec:troubleshoot-couvertures-perte-mise-a-jour}

Vous avez mis à jour depuis une install antérieure à 1.1.4 et toutes
les couvertures précédemment cataloguées ont disparu. Il s'agit de
l'issue~\#213 — l'ancien modèle compose ne persistait pas
\texttt{/app/covers}, donc le layer en écriture portant ces JPG a
été détruit quand le nouveau conteneur a remplacé l'ancien.

Correctif : récupérez le \texttt{docker-compose.yml} de la 1.1.4,
puis lancez \texttt{docker compose down \&\& docker compose up -d}.
Les couvertures suivantes seront persistées dans le volume
\texttt{mybibli\_covers}.

Récupération des couvertures perdues : ouvrez la page de détail de
chaque titre concerné et cliquez sur \emph{Re-télécharger les
métadonnées} pour repeupler depuis la chaîne de fournisseurs. Une
action en masse est suivie dans l'issue~\#214.

\section{Où signaler un bug}

Problèmes que vous ne pouvez pas résoudre, ou comportement qui
contredit ce manuel : ouvrez un ticket sur
\url{https://github.com/guycorbaz/mybibli/issues}. Le template
d'issue demande la version (visible en pied de page), la méthode
d'installation et les étapes pour reproduire — merci de les
renseigner.
